\section{Change of prescribed base seminorm}\label{sec:general-base-change}

Let $R\to S$ be a homomorphism of commutative rings and let $v$ on $R$
and $w$ on $S$ be real-valued multiplicative ultrametric seminorms.
Their compatibility means $w(r)=v(r)$ for every coefficient from $R$.
All compatibility conditions in this section are explicit hypotheses.

\begin{definition}[Restriction of coefficient seminorms]
\label{spectrum:scalar-restriction}
\lean{EnvelopeProperty.Berkovich.RelativeSpectrum.restrictScalars,
EnvelopeProperty.Berkovich.RelativeSpectrum.restrictScalars_val_apply,
EnvelopeProperty.Berkovich.RelativeSpectrum.continuous_restrictScalars,
EnvelopeProperty.Berkovich.RelativeSpectrum.restrictScalars_comap}
\leanok
\uses{spectrum:relative}
For an $S$-algebra $B$ with its compatible $R$-algebra structure,
every point of $\mathcal S_w(B)$ is a point of $\mathcal S_v(B)$
by restriction of coefficients. This map leaves its entire seminorm
on $B$ unchanged and is continuous.
\end{definition}
\begin{proof}
For $r\in R$, its image in $B$ is the image of its image in $S$.
The $B$-seminorm therefore takes value $w(r)=v(r)$ there.
The resulting absolute spectrum point is unchanged, so all evaluation
functions are unchanged and continuity follows from the subspace topology.
\end{proof}

\begin{theorem}[Restriction after scalar extension]
\label{spectrum:base-change-restriction}
\lean{EnvelopeProperty.Berkovich.RelativeSpectrum.baseChangeRestriction,
EnvelopeProperty.Berkovich.RelativeSpectrum.baseChangeRestriction_val_apply,
EnvelopeProperty.Berkovich.RelativeSpectrum.continuous_baseChangeRestriction,
EnvelopeProperty.Berkovich.RelativeSpectrum.baseChange_val_tmul,
EnvelopeProperty.Berkovich.RelativeSpectrum.baseChangeRestriction_val_aeval}
\leanok
\uses{spectrum:scalar-restriction,spectrum:pullback}
For any commutative $R$-algebra $A$, restriction along
$A\to S\otimes_R A$, $a\mapsto1\otimes a$, defines a continuous map
\[
 \rho:\mathcal S_w(S\otimes_R A)\longrightarrow\mathcal S_v(A).
\]
For every $a\in A$ and $s\in S$,
$|a|_{\rho(x)}=|1\otimes a|_x$ and
$|s\otimes a|_x=w(s)|a|_{\rho(x)}$.
For every multivariate polynomial $p\in R[U_i]$ and family $a_i\in A$,
the value of $p(a_i)$ under $\rho(x)$ equals the value of
$p(1\otimes a_i)$ under $x$.
\end{theorem}
\begin{proof}
Restrict coefficients from $S$ to $R$, then pull back along the actual
tensor inclusion. This proves the value formula for every $a$ and
continuity. A pure tensor is $(s\otimes1)(1\otimes a)$.
The first factor is a base coefficient, so multiplicativity gives the
pure-tensor formula. The tensor inclusion is an algebra homomorphism
on the whole ring, and therefore commutes with polynomial evaluation,
proving the last assertion on every polynomial.
\end{proof}

\begin{theorem}[Algebra-map functoriality]
\label{spectrum:base-change-algebra-naturality}
\lean{EnvelopeProperty.Berkovich.RelativeSpectrum.baseChangeMap,
EnvelopeProperty.Berkovich.RelativeSpectrum.baseChangeMap_tmul,
EnvelopeProperty.Berkovich.RelativeSpectrum.baseChangeRestriction_comap,
EnvelopeProperty.Berkovich.RelativeSpectrum.baseChangeMap_id,
EnvelopeProperty.Berkovich.RelativeSpectrum.baseChangeMap_comp}
\leanok
\uses{spectrum:base-change-restriction}
For an $R$-algebra homomorphism $\varphi:A\to B$, the actual tensor
homomorphism $\mathrm{id}_S\otimes\varphi$ induces continuous pullback
on the $w$-relative spectra. Restriction commutes with this pullback
and with $\varphi^*$ on every ring element. Tensor pullbacks preserve
identities and reverse the composition of algebra homomorphisms.
\end{theorem}
\begin{proof}
The tensor homomorphism sends $s\otimes a$ to $s\otimes\varphi(a)$
and preserves every $S$ coefficient. Its composition with the right
inclusion equals the right inclusion after $\varphi$, as homomorphisms
on the full algebra. Pulling back gives the restriction square.
Identity and composition follow from the corresponding identities
of the actual tensor algebra homomorphisms.
\end{proof}

\begin{theorem}[Compatible scalar-map functoriality]
\label{spectrum:base-change-scalar-naturality}
\lean{EnvelopeProperty.Berkovich.RelativeSpectrum.scalarChangeComap,
EnvelopeProperty.Berkovich.RelativeSpectrum.scalarChangeComap_val_apply,
EnvelopeProperty.Berkovich.RelativeSpectrum.continuous_scalarChangeComap,
EnvelopeProperty.Berkovich.RelativeSpectrum.baseChangeRestriction_scalarChangeComap,
EnvelopeProperty.Berkovich.RelativeSpectrum.scalarChangeComap_id,
EnvelopeProperty.Berkovich.RelativeSpectrum.scalarChangeComap_comp}
\leanok
\uses{spectrum:base-change-restriction}
Let $\tau:S\to T$ be an $R$-algebra homomorphism and let $u$ on $T$
satisfy $u\circ\tau=w$. Pullback along
$\tau\otimes\mathrm{id}_A$ defines a continuous map
$\mathcal S_u(T\otimes_R A)\to\mathcal S_w(S\otimes_R A)$.
It commutes with restriction to $A$. These maps obey identity and
composition laws for compatible scalar maps.
\end{theorem}
\begin{proof}
The tensor homomorphism sends an $S$ coefficient $s\otimes1$ to
$\tau(s)\otimes1$, whose seminorm is $u(\tau(s))=w(s)$.
It sends $1\otimes a$ to $1\otimes a$; this proves the restriction
square on every $a\in A$. Continuity and the functorial identities
follow from full homomorphism pullback and tensor map identities.
\end{proof}

\begin{theorem}[Direct and iterated scalar extension]
\label{spectrum:base-change-iteration}
\lean{EnvelopeProperty.Berkovich.RelativeSpectrum.scalarCompatibility_tower,
EnvelopeProperty.Berkovich.RelativeSpectrum.iteratedBaseChangeHomeomorph,
EnvelopeProperty.Berkovich.RelativeSpectrum.iteratedBaseChangeHomeomorph_val_tmul,
EnvelopeProperty.Berkovich.RelativeSpectrum.iteratedBaseChangeHomeomorph_symm_val_tmul,
EnvelopeProperty.Berkovich.RelativeSpectrum.baseChangeRestriction_iterated}
\leanok
\uses{spectrum:base-change-restriction}
For a compatible tower $R\to S\to T$, cancellation of the middle
tensor factor gives a canonical homeomorphism
\[
 \mathcal S_u\bigl(T\otimes_S(S\otimes_R A)\bigr)
 \simeq\mathcal S_u(T\otimes_R A).
\]
It identifies successive restrictions to $A$ with direct restriction.
\end{theorem}
\begin{proof}
The actual tensor algebra isomorphism sends
$t\otimes(s\otimes a)$ to $st\otimes a$, with $s$ mapped into $T$.
Its inverse sends $t\otimes a$ to $t\otimes(1\otimes a)$.
Pullback along these inverse algebra homomorphisms gives inverse
continuous maps of the spectra with prescribed base seminorm $u$.
Under this comparison, the representative of every original $a\in A$
is $1\otimes(1\otimes a)$ in both routes. Hence the two restrictions
agree on every original ring element and thus as entire spectrum points.
The compatibility of $u$ with $v$ follows by composing its explicit
compatibility with $w$ and that of $w$ with $v$.
\end{proof}

These are genuine restriction maps. No assertion that base-change
restriction is surjective, and no seminorm-extension existence
hypothesis, is part of these constructions.
